\documentclass{bmstu}

\usepackage{svg}
\usepackage{adjustbox}

\usetikzlibrary{positioning,shapes.geometric}

\addbibresource{references.bib}

% =============================================================================
% ШАБЛОН ОТЧЁТА ПО ЛР №1 «Сравнительный анализ ресурсной эффективности алгоритмов»
% -----------------------------------------------------------------------------
% Как пользоваться шаблоном:
%  - Каждая строка, начинающаяся с «%», является ИНСТРУКЦИЕЙ. Под ней (или
%    вместо неё) нужно написать собственный текст. Перед сдачей инструкции
%    можно удалить: на результат компиляции они не влияют.
%  - Заголовки глав, разделов и подразделов НЕ удалять, НЕ переименовывать
%    и НЕ менять местами: план отчёта фиксирован.
%  - Задача по варианту у всех одинакова: поиск элемента x в массиве A из N
%    целых чисел. Требуется четыре алгоритма:
%      1) линейный поиск;
%      2) бинарный поиск, нерекурсивная реализация, вариант 1
%         (с досрочным выходом при совпадении);
%      3) бинарный поиск, нерекурсивная реализация, вариант 2
%         (без досрочного выхода, проверка на равенство один раз после цикла);
%      4) бинарный поиск, рекурсивная реализация.
%  - Обозначения во всём отчёте должны быть одинаковыми: A — массив, N — его
%    размер, x — искомое значение, l и r — границы диапазона, mid — индекс
%    середины. Если выбрали другие обозначения, используйте их везде.
%  - Формулы оформляются в математическом режиме: $A_{mid}$, $O(\log_2 N)$.
%    Индекс из нескольких букв обязательно брать в фигурные скобки:
%    $A_{mid}$, а не $A_mid$.
%  - На каждый рисунок, таблицу и листинг должна быть ссылка в тексте
%    (\ref{...}), причём ссылка должна стоять ДО самого объекта.
%  - Стиль изложения безличный: «разработан», «выполнено», «приведено».
%    Не пишите «я сделал», «мы видим».
%  - Длинное тире в тексте: «---» (три дефиса). Маркер списка «---»
%    задаётся командой \renewcommand{\labelitemi}{---} во введении.
% =============================================================================

\begin{document}

\makereporttitle
{Информатика и системы управления}
{Программное обеспечение ЭВМ и информационные технологии}
{лабораторной работе №1}
{Анализ алгоритмов}
{Сравнительный анализ ресурсной эффективности алгоритмов}
{} % Вариант
{М. А. Кузнецов/ИУ7-52Б} % Номер группы/ФИО студента (если несколько - разделить запятой)
{Д. В. Строганов}
% ИНСТРУКЦИЯ: замените ФИО и группу студента на свои. Больше на титульном
% листе ничего не менять.

\maketableofcontents

\chapter*{ВВЕДЕНИЕ}
\addcontentsline{toc}{chapter}{ВВЕДЕНИЕ}

% АБЗАЦ 1. Цель работы — одно предложение. Формулировку взять ДОСЛОВНО из
% задания преподавателя. Рекомендуемая форма:
%   «Цель работы --- выполнение сравнительного анализа ресурсной
%    эффективности алгоритмов.»

% АБЗАЦ 2. Фраза «Задачи работы:» и следом маркированный список.
% Перед списком поставить \renewcommand{\labelitemi}{---}, сам список
% оформить окружением itemize. Пункты списка переносятся ДОСЛОВНО из
% задания преподавателя, в том же порядке, каждый с маленькой буквы
% и с «;» в конце (последний пункт с «.»). Должно быть ровно 11 пунктов:
%   описать решаемую задачу согласно варианту;
%   разработать рекурсивный и нерекурсивные алгоритмы решения задачи,
%     согласно варианту;
%   описать разработанные алгоритмы схемами алгоритмов;
%   описать требования к программному обеспечению;
%   описать средства реализации алгоритмов;
%   реализовать разработанные алгоритмы;
%   выполнить тестирование реализации алгоритмов;
%   выполнить теоретическую оценку затрачиваемой реализацией каждого
%     алгоритма памяти (для рекурсивного алгоритма --- на материале анализа
%     высоты дерева рекурсивных вызовов и оценки затрачиваемой на один
%     вызов функции памяти);
%   оценить временную эффективность алгоритмов или их реализаций в лучшем
%     и худшем случае в терминах числа сравнений искомой величины
%     с другими величинами;
%   представить результаты сравнительного анализа в формате таблицы;
%   сделать выводы из сравнительного анализа реализации разработанных
%     алгоритмов решения одной и той же задачи, заданной вариантом, по
%     критериям временной и ёмкостной эффективности, а также дополнительных
%     требований, налагаемых алгоритмом.
% Проверьте, что ни у одного пункта не пропущена «;».

\chapter{Аналитическая часть}

\section{Постановка задачи}
% АБЗАЦ 1. Формальная постановка задачи в 1–2 предложениях. Обязательно:
%  - что дано: массив A из N целых чисел и искомое значение x;
%  - что требуется: определить, содержится ли x в A, и если да,
%    вернуть его индекс.
% Пример формулировки: «Дан массив $A$ из $N$ целых чисел. Для заданного
% значения $x$ требуется определить, содержится ли $x$ в массиве $A$,
% и если содержится --- вернуть его индекс в массиве.»

% АБЗАЦ 2. Фраза «Требуется разработать следующие алгоритмы:» и
% маркированный список (itemize) из четырёх пунктов:
%   линейный поиск;
%   бинарный поиск, нерекурсивная реализация, вариант 1;
%   бинарный поиск, нерекурсивная реализация, вариант 2;
%   бинарный поиск, рекурсивная реализация.
% Названия алгоритмов здесь должны в точности совпадать с названиями
% во всех последующих разделах, подписях рисунков и в итоговой таблице.

% АБЗАЦ 3 (по желанию, одно предложение). Указать, что предметом анализа
% является сравнение этих реализаций между собой по времени и памяти.

% АБЗАЦ 4. Входные и выходные данные, явно и раздельно:
%  - входные данные: натуральное число N, массив из N целых чисел,
%    искомое целое значение x;
%  - выходные данные: индекс элемента, равного x, либо $-1$, если такого
%    элемента в массиве нет.
% Можно оформить двумя предложениями: «Входные данные: ...
% Выходные данные: ...».

\chapter{Конструкторская часть}

\section{Разработка алгоритмов}
% АБЗАЦ 1 (вводный, одно-два предложения). Перечислить, какие алгоритмы
% разработаны (те же четыре, что в постановке задачи), и сослаться на
% источник, где описаны линейный и бинарный поиск: ~\cite{ключ_источника}.
% Источник должен быть добавлен в файл references.bib (например, Кнут
% «Искусство программирования», т. 3, или Скиена «The Algorithm Design
% Manual»). Ключ в \cite должен совпадать с ключом в references.bib.
%
% (По желанию) Отдельным абзацем ввести общие обозначения, которые
% используются на схемах: A — массив, N — размер, x — искомое значение,
% l, r — левая и правая границы диапазона, mid — индекс середины,
% res/index — результат. Это облегчает чтение схем.

\subsubsection{Линейный алгоритм}
% АБЗАЦ 1. Словесное описание линейного поиска, 3–4 предложения.
% Обязательно упомянуть:
%  - что алгоритм нерекурсивный;
%  - что массив НЕ обязан быть отсортирован;
%  - что элементы $A_i$ перебираются последовательно от индекса 0 до N-1,
%    и каждый сравнивается с x;
%  - что при совпадении возвращается индекс i и работа завершается;
%  - что если совпадений нет, возвращается $-1$.
% Последнее предложение абзаца: «Схема алгоритма приведена на
% рисунке~\ref{fig:linear-search}.»

% РИСУНОК. Схема алгоритма линейного поиска (ГОСТ 19.701-90: блоки
% «Начало/Конец», «Процесс», «Решение», «Граница цикла»). Вставить через
% окружение figure[H] + \includesvg (или \includegraphics для PNG/PDF),
% обёрнутое в adjustbox с max width=\textwidth, max height=0.8\textheight.
% Подпись: \caption{Схема линейного поиска}
% Метка:   \label{fig:linear-search}
% На схеме должны быть видны: инициализация результата, цикл по i,
% условие $A_i = x$, ветка возврата индекса, возврат $-1$.

\subsubsection{Нерекурсивный алгоритм бинарного поиска}
% АБЗАЦ 1. Общий принцип нерекурсивного бинарного поиска, 3–5 предложений.
% Обязательно упомянуть:
%  - что алгоритм работает только на массиве, ОТСОРТИРОВАННОМ по
%    возрастанию (в программе — на отсортированной копии исходного массива);
%  - что поиск ведётся в диапазоне между индексами l и r, изначально
%    l = 0, r = N - 1 (границы всего массива);
%  - что на каждой итерации вычисляется середина диапазона
%    mid = l + \lfloor (r - l)/2 \rfloor (можно вынести в отдельную формулу);
%  - что элемент $A_{mid}$ сравнивается с x и диапазон сужается вдвое:
%    если $A_{mid} < x$, то поиск продолжается в правой половине
%    (l = mid + 1), иначе — в левой (r = mid или r = mid - 1, в зависимости
%    от варианта).
% ВНИМАНИЕ: границы сдвигаются на ИНДЕКС mid, а не на значение $A_{mid}$.
% Проверьте, что направление сдвига указано верно: при $A_{mid} < x$
% двигается ЛЕВАЯ граница l.

% АБЗАЦ 2. Различия двух вариантов. Описать каждый вариант 2–3
% предложениями:
%  - Вариант 1: на каждой итерации сначала проверяется равенство
%    $A_{mid} = x$; при совпадении работа завершается досрочно и
%    возвращается mid. Цикл идёт, пока l <= r (замкнутый диапазон). Если
%    диапазон исчерпан (стал пустым), возвращается $-1$.
%  - Вариант 2: на равенство на каждой итерации НЕ проверяется, выполняется
%    только одно сравнение $A_{mid} < x$. Цикл идёт, пока l < r, то есть
%    пока диапазон не сократится до одного элемента. После цикла один раз
%    проверяется $A_l = x$: при равенстве возвращается l, иначе $-1$.
% Полезно добавить фразу о следствии различия: вариант 2 не завершается
% досрочно, зато выполняет меньше сравнений на одной итерации; при
% нескольких одинаковых элементах вариант 2 находит первое вхождение.
% Описание обязательно должно совпадать с тем, что изображено на схемах
% и реализовано в коде.
% Последнее предложение абзаца: «Схемы алгоритмов приведены на
% рисунке~\ref{fig:binary-search-iterative-v1} и
% рисунке~\ref{fig:binary-search-iterative-v2}.»

% РИСУНОК. Схема нерекурсивного бинарного поиска, вариант 1.
% Оформление как у предыдущего рисунка.
% Подпись: \caption{Схема нерекурсивного бинарного поиска, вариант 1}
% Метка:   \label{fig:binary-search-iterative-v1}
% На схеме должны быть: инициализация l, r (и переменной результата),
% цикл с условием l <= r, вычисление mid, проверка $A_{mid} = x$ с выходом,
% проверка $A_{mid} < x$ с двумя ветками сдвига границ, возврат результата.

% РИСУНОК. Схема нерекурсивного бинарного поиска, вариант 2.
% Подпись: \caption{Схема нерекурсивного бинарного поиска, вариант 2}
% Метка:   \label{fig:binary-search-iterative-v2}
% На схеме должны быть: инициализация l, r, цикл с условием l < r,
% вычисление mid, ОДНА проверка $A_{mid} < x$ с двумя ветками
% (l = mid + 1 / r = mid), после цикла проверка $A_l = x$ и возврат
% l или $-1$.

\subsubsection{Рекурсивный алгоритм бинарного поиска}
% АБЗАЦ 1. Описание рекурсивного бинарного поиска, 4–5 предложений.
% Обязательно упомянуть:
%  - что используется тот же принцип деления диапазона пополам, что и
%    в нерекурсивных вариантах;
%  - что повторение организовано НЕ циклом, а рекурсивным вызовом функции
%    для суженного диапазона: новые границы l и r передаются как аргументы
%    (в первом вызове l = 0, r = N - 1);
%  - условие выхода из рекурсии (базовый случай): например, l > r —
%    элемент не найден, вернуть $-1$; и/или $A_{mid} = x$ — вернуть mid.
%    Описать ровно тот вариант, который реализован в коде (с досрочным
%    выходом, как вариант 1, или без него, как вариант 2) — от этого
%    зависят оценки в исследовательской части;
%  - (по желанию) если в коде есть функция-обёртка, которая вызывает
%    рекурсивную функцию с начальными границами, упомянуть её.
% Отдельным предложением подчеркнуть принципиальное отличие от
% нерекурсивных реализаций: самовызов функции вместо цикла. Именно это
% определяет разницу в затрачиваемой памяти (сослаться на раздел оценки
% памяти).
% Последнее предложение абзаца: «Схема алгоритма представлена на
% рисунке~\ref{fig:binary-search-recursive}.»

% РИСУНОК. Схема рекурсивного бинарного поиска.
% Подпись: \caption{Схема рекурсивного бинарного поиска}
% Метка:   \label{fig:binary-search-recursive}
% На схеме должны быть: заголовок функции с параметрами (A, x, l, r),
% проверка базового случая, вычисление mid, сравнения и ДВА блока
% рекурсивного вызова (для левой и правой половины), возврат результата.

\section{Требования к программному обеспечению}
% АБЗАЦ 1. Фраза «Программное обеспечение должно:» и маркированный список.
% Требования взяты с семинара 1 и у всех одинаковы. Обязательные пункты:
%  - запрашивать у пользователя размер массива N, затем запрашивать
%    с клавиатуры N целых чисел, формируя массив A;
%  - создавать копию массива A, сортировать её по возрастанию и выводить
%    отсортированный массив на экран (исходный массив не изменяется;
%    линейный поиск работает по исходному, бинарные — по копии);
%  - циклически выводить меню, через которое пользователь вводит искомое
%    значение x и выбирает алгоритм поиска;
%  - быть реализовано на языке программирования со статической
%    типизацией.
% (По желанию) можно перечислить пункты меню: выбор алгоритма, поиск,
% поиск всеми алгоритмами, вывод массива, запуск тестов, выход.
% Каждый пункт начинается с глагола в неопределённой форме,
% с маленькой буквы, заканчивается «;», последний — «.».

\chapter{Технологическая часть}

\section{Средства реализации}
% АБЗАЦ 1. Выбранный язык программирования, 1–2 предложения.
% Обязательно:
%  - название языка и стандарт/версию (например, C++20, C++17, Swift 5);
%  - обоснование выбора: язык со статической типизацией, как требуют
%    требования к ПО (можно добавить ещё одну причину: наличие нужных
%    контейнеров в стандартной библиотеке, знакомство с языком и т.п.);
%  - ссылку на источник со стандартом или документацией языка:
%    \cite{ключ_источника} (добавить в references.bib).

% АБЗАЦ 2. Фраза «Для реализации использованы средства стандартной
% библиотеки:» и маркированный список конкретных средств с пояснением
% через тире, для чего каждое нужно. Например:
%  - \texttt{std::vector} --- динамический контейнер для хранения массива;
%  - средства ввода-вывода (iostream или scanf/printf);
%  - \texttt{std::sort} --- сортировка копии массива (если используется).
% Названия функций и типов оформлять через \texttt{...}.

% АБЗАЦ 3. Компилятор/среда сборки и флаги, одно предложение.
% Например: «Для сборки использован компилятор GCC с флагами
% \texttt{-std=c++20 -Wall -Wextra -O0}.» Указывать реальные флаги из
% вашего Makefile/скрипта сборки.

\section{Реализация алгоритмов}
% АБЗАЦ 1. Описание структуры исходного кода, 2–3 предложения. Для каждого
% файла: имя файла (\texttt{...}), что в нём находится, и ссылка на
% листинг в приложении (Листинг~\ref{lst:...}). Например:
%  - заголовочный файл с объявлениями функций поиска;
%  - файл с определениями (реализациями) функций поиска;
%  - файл с главной функцией и интерактивным меню.
% Метки листингов должны совпадать с теми, что заданы в Приложении А.
% Сам код в этот раздел не вставлять — он целиком лежит в приложении.
% (Допустимая альтернатива: вставить листинги функций поиска прямо сюда,
% окружением lstlisting/lstinputlisting, с подписью вида «Линейный поиск»;
% тогда в приложении оставить только main и тесты.)

\section{Тестирование реализации алгоритмов}
% АБЗАЦ 1. Как проводилось тестирование, 2–3 предложения. Обязательно:
%  - имя файла модуля тестирования и ссылка на листинг
%    (Листинг~\ref{lst:test-cpp});
%  - что одни и те же тесты прогоняются для ВСЕХ четырёх реализаций;
%  - перечень классов тестовых случаев (одним предложением через «;»):
%    поиск элемента в начале, в середине и в конце массива; поиск
%    отсутствующего элемента: меньше минимального, больше максимального,
%    внутри диапазона значений; поиск в массиве из одного элемента
%    (элемент есть / элемента нет).
% (По желанию, это усиливает отчёт) добавить: массив из двух элементов,
% массивы чётной и нечётной длины, отрицательные числа.

% АБЗАЦ 2. Одно предложение с тестовым массивом, на котором выполнено
% большинство тестов, например: «Для тестирования использован массив
% $\{...\}$.» Массив должен быть отсортирован, длиной 8–10 элементов,
% с отрицательными и положительными числами.
% Перед таблицей добавить ссылку: «Описание тестов приведено в
% таблице~\ref{tab:tests}.»

% ТАБЛИЦА. Окружение table[H] + tabularx на всю ширину (\textwidth).
% Подпись: \caption{Описание тестов}   Метка: \label{tab:tests}
% (подпись ставится НАД таблицей, т.е. \caption до tabularx).
% Столбцы (можно добавить первым столбец «№»):
%   Описание теста | Входные данные | Ожидаемый результат | Полученный результат
% Строки — по одной на каждый тестовый случай из абзаца 1 (минимум 8):
%   - элемент в начале массива          -> индекс 0;
%   - элемент в конце массива           -> индекс N-1;
%   - элемент в середине массива        -> его индекс;
%   - элемент меньше минимального       -> «Не найден ($-1$)»;
%   - элемент больше максимального      -> «Не найден ($-1$)»;
%   - отсутствующий элемент внутри диапазона -> «Не найден ($-1$)»;
%   - массив из одного элемента, элемент есть  -> индекс 0;
%   - массив из одного элемента, элемента нет  -> «Не найден ($-1$)».
% Во «Входных данных» указать x (и массив, если он отличается от
% основного). «Полученный результат» заполнять ТОЛЬКО по факту реального
% запуска тестов, а не переписывать ожидаемый.

% АБЗАЦ 3. Итог тестирования, одно предложение: «Все тесты пройдены
% успешно.» (Если какой-то тест не прошёл — исправить код и перезапустить.)

\chapter{Исследовательская часть}

\section{Оценка затрачиваемой памяти}
% АБЗАЦ 1 (вводный, 1–2 предложения). Из чего складывается память,
% затрачиваемая реализацией: локальные (служебные) переменные функции,
% а для рекурсивной реализации — кадры стека вызовов. Уточнить, что
% оценивается ДОПОЛНИТЕЛЬНАЯ память, т.е. без учёта самого входного
% массива (массив передаётся по ссылке и не копируется).

% АБЗАЦ 2. Линейный поиск, 2 предложения. Перечислить конкретные
% служебные переменные реализации (например, счётчик i, результат);
% указать, что их число фиксировано и не зависит от N. Вывод:
% дополнительная память $O(1)$.

% АБЗАЦ 3. Нерекурсивный бинарный поиск (оба варианта), 2–3 предложения.
% Перечислить переменные: l, r, mid (и переменная-результат, если есть).
% Число переменных фиксировано, не зависит от N, следовательно
% дополнительная память $O(1)$. Отдельно отметить, что варианты 1 и 2
% используют одинаковый набор переменных, поэтому затрачивают одинаковую
% память.

% АБЗАЦ 4. Рекурсивный бинарный поиск — ГЛАВНАЯ часть раздела, этого
% прямо требует задание. По шагам:
%  (а) Что хранится в ОДНОМ кадре стека: аргументы функции (ссылка/
%      указатель на массив, x, l, r), локальная переменная mid, адрес
%      возврата, сохранённые регистры/служебная информация компилятора.
%      Сказать, что размер кадра постоянен, обозначить его c.
%      (По желанию, но очень желательно) оценить c в байтах для своей
%      платформы: например, указатель — 8 байт, int — 4 байта, адрес
%      возврата — 8 байт и т.д.; итог записать формулой.
%  (б) Кадр существует до завершения своего вызова, поэтому максимальная
%      память определяется высотой дерева рекурсивных вызовов h
%      (максимальной глубиной рекурсии).
%  (в) Дерево вызовов вырождено в цепочку (каждый вызов порождает не
%      более одного рекурсивного вызова), диапазон на каждом вызове
%      уменьшается вдвое, поэтому в худшем случае
%        $$h = \lceil \log_2 (N + 1) \rceil$$
%      (для реализации с досрочным выходом и условием l > r).
%      Если ваша реализация без досрочного выхода (как вариант 2),
%      пересчитайте h для неё, например $h = \lceil \log_2 N \rceil + 1$.
%      Формула должна соответствовать вашему коду!
%  (г) Итоговая формула максимальной затрачиваемой памяти:
%        $$M_{\max}(N) = c \cdot h = c \cdot \lceil\log_2(N + 1)\rceil
%          = O(\log_2 N).$$
%  (д) Заключительное предложение: рекурсивная реализация требует
%      $O(\log_2 N)$ дополнительной памяти против $O(1)$ у нерекурсивных.
% Знак препинания после выключной формулы ставить ВНУТРИ $$...$$.

\section{Оценка временной эффективности}
% АБЗАЦ 1 (вводный, 1–2 предложения). Временная эффективность оценивается
% числом сравнений искомого значения x с элементами массива (именно
% так сказано в задании). Сравнения индексов (l <= r, i < N) не
% учитываются — это стоит явно оговорить.
% Ввести обозначения: $T_{\min}(N)$ — число сравнений в лучшем случае,
% $T_{\max}(N)$ — в худшем случае.
%
% Во всех подразделах ниже для КАЖДОГО алгоритма нужно:
%  1) описать, какая ситуация является лучшим случаем, и выписать
%     $T_{\min}(N)$ точной формулой и асимптотикой;
%  2) описать, какая ситуация является худшим случаем, и выписать
%     $T_{\max}(N)$ точной формулой и асимптотикой.
% Можно оформлять текстом или списком «лучший случай: ...; худший
% случай: ...».

\subsection{Линейный поиск}
% АБЗАЦ 1.
%  - Лучший случай: x находится на первой позиции массива — 1 сравнение,
%    $T_{\min}(N) = 1 = O(1)$.
%  - Худший случай: x на последней позиции или отсутствует — проверяются
%    все элементы, N сравнений, $T_{\max}(N) = N = O(N)$.

\subsection{Бинарный поиск, нерекурсивная реализация}
% АБЗАЦ 1 (вводный, по желанию). Пояснить, сколько сравнений с x делается
% на одной итерации в каждом варианте:
%  - вариант 1: итерация без совпадения — 2 сравнения (на равенство
%    и на порядок), итерация с совпадением — 1 сравнение;
%  - вариант 2: на каждой итерации 1 сравнение, плюс одна проверка на
%    равенство после цикла.
% (По желанию) показать вывод числа итераций: размер диапазона после
% i-й итерации $N_i = N / 2^i$, поиск заканчивается при $N_i = 1$,
% откуда $i \approx \log_2 N$.

% АБЗАЦ 2. Вариант 1:
%  - Лучший случай: x совпадает с центральным элементом исходного
%    диапазона — 1 сравнение, $T_{\min}(N) = 1 = O(1)$.
%  - Худший случай: x отсутствует (или найден на последнем шаге) —
%    $\lceil\log_2(N+1)\rceil$ итераций по 2 сравнения,
%    $T_{\max}(N) = 2\lceil\log_2(N+1)\rceil = O(\log_2 N)$.

% АБЗАЦ 3. Вариант 2:
%  - Досрочного выхода нет, поэтому лучший и худший случаи отличаются
%    только числом итераций (на 1), но не асимптотикой.
%  - Лучший случай: $\lfloor\log_2 N\rfloor$ итераций и одна итоговая
%    проверка, $T_{\min}(N) = \lfloor\log_2 N\rfloor + 1 = O(\log_2 N)$.
%  - Худший случай: $\lceil\log_2 N\rceil$ итераций и одна итоговая
%    проверка, $T_{\max}(N) = \lceil\log_2 N\rceil + 1 = O(\log_2 N)$.
% Завершающее предложение: лучший случай варианта 2 асимптотически хуже,
% чем у варианта 1 ($O(\log_2 N)$ против $O(1)$), а худший случай
% асимптотически совпадает.

\subsection{Бинарный поиск, рекурсивная реализация}
% АБЗАЦ 1. Оценка зависит от того, как устроена ВАША рекурсивная
% функция, и должна совпадать с описанием в конструкторской части:
%  - если рекурсия с досрочным выходом (проверка $A_{mid} = x$ на каждом
%    вызове): оценки как у варианта 1 —
%    $T_{\min}(N) = 1 = O(1)$ (x в центре исходного диапазона),
%    $T_{\max}(N) = 2\lceil\log_2(N+1)\rceil = O(\log_2 N)$
%    ($\lceil\log_2(N+1)\rceil$ рекурсивных вызовов по 2 сравнения);
%  - если без досрочного выхода: оценки как у варианта 2.
% Добавить предложение, что число рекурсивных шагов в худшем случае
% совпадает с высотой дерева рекурсивных вызовов h из раздела оценки
% памяти: по числу сравнений рекурсия не отличается от соответствующего
% нерекурсивного варианта, отличие только в памяти.

\section{Результаты сравнительного анализа}
% АБЗАЦ 1. Одно предложение со ссылкой на таблицу: «Результаты
% теоретической оценки ресурсной эффективности алгоритмов приведены
% в таблице~\ref{tab:comparison}.»

% ТАБЛИЦА. Окружение table[H] + tabularx на всю ширину.
% Подпись: \caption{Сравнительный анализ ресурсной эффективности алгоритмов}
% Метка:   \label{tab:comparison}
% Столбцы:
%   Реализация | $T_{\min}(N)$ | $T_{\max}(N)$ |
%   Требования, накладываемые алгоритмом | Дополнительная память
% Строки (четыре, в том же порядке, что в постановке задачи), значения
% должны СОВПАДАТЬ с полученными выше в разделах оценки:
%   Линейный поиск:                     O(1)      | O(N)      | нет требований (массив может быть не отсортирован) | O(1)
%   Бинарный, нерекурсивный, вариант 1: O(1)      | O(log2 N) | массив предварительно отсортирован                  | O(1)
%   Бинарный, нерекурсивный, вариант 2: O(log2 N) | O(log2 N) | массив предварительно отсортирован                  | O(1)
%   Бинарный, рекурсивный:              по вашей реализации | O(log2 N) | массив отсортирован; ограничение глубины рекурсии размером стека | O(log2 N)
% (По желанию) можно добавить столбцы с точными формулами числа
% сравнений, а не только асимптотикой.

% АБЗАЦ 2 (по желанию, но желательно). Пояснение к таблице про
% дополнительное требование рекурсивной реализации: глубина рекурсии
% ограничена размером стека вызовов, при очень больших N теоретически
% возможно переполнение стека; у нерекурсивных реализаций такого
% ограничения нет. Также можно отметить, что требование
% отсортированности означает затраты на предварительную сортировку
% (обычно $O(N\log N)$), которые окупаются при многократном поиске.

\section{Выводы}
% АБЗАЦ 1. Фраза «По результатам сравнительного анализа сделаны следующие
% выводы:» и маркированный список из 4–5 пунктов. Каждый вывод должен
% опираться на конкретные значения из таблицы и ЗАКАНЧИВАТЬСЯ сравнением
% «... против ...». Обязательно охватить три критерия задания: время,
% память, дополнительные требования. Ориентиры содержания:
%  - бинарный поиск асимптотически эффективнее линейного по времени
%    в худшем случае ($O(\log_2 N)$ против $O(N)$), но требует, чтобы
%    массив был предварительно отсортирован;
%  - вариант 1 нерекурсивного бинарного поиска и рекурсивная реализация
%    асимптотически эквивалентны по времени (если рекурсия с досрочным
%    выходом), но нерекурсивная реализация требует меньше памяти
%    ($O(1)$ против $O(\log_2 N)$);
%  - варианты 1 и 2 нерекурсивного бинарного поиска различаются только
%    лучшим случаем: $O(1)$ против $O(\log_2 N)$; память у них одинакова.
%    Вариант 2 полезен, когда важна предсказуемость числа сравнений
%    или нужно найти первое вхождение;
%  - рекурсивная реализация наихудшая по памяти и дополнительно
%    ограничена размером стека вызовов;
%  - итоговая рекомендация: какая реализация предпочтительна для решения
%    задачи и почему (как правило, нерекурсивный бинарный поиск,
%    вариант 1, если массив отсортирован; линейный — если сортировать
%    невыгодно, например при однократном поиске).
% Проверьте, что асимптотики в выводах совпадают с таблицей.

\chapter*{ЗАКЛЮЧЕНИЕ}
\addcontentsline{toc}{chapter}{ЗАКЛЮЧЕНИЕ}
% АБЗАЦ 1. Утверждение о достижении цели, повторяющее формулировку цели
% из введения: «Цель работы достигнута: выполнен сравнительный анализ
% ресурсной эффективности алгоритмов.»
% Далее 2–3 предложения с главными итогами (кратко, без формул вывода):
% какая реализация лучшая по времени и памяти, какая худшая по памяти,
% какая худшая по времени. Итоги не должны противоречить разделу
% «Выводы».

% АБЗАЦ 2. Фраза «В ходе работы были решены поставленные задачи:» и
% маркированный список. Это те же 11 задач из введения в том же порядке,
% но переписанные в прошедшем времени, страдательном залоге:
%   «описать решаемую задачу» -> «описана решаемая задача согласно варианту»;
%   «разработать алгоритмы»   -> «разработаны рекурсивный и нерекурсивные
%                                  алгоритмы согласно варианту»;
%   «выполнить тестирование»  -> «выполнено тестирование реализации
%                                  алгоритмов»; и т.д.
% Должно быть ровно 11 пунктов, по одному на каждую задачу; ни одну
% задачу не пропускать. Пункты разделяются «;», последний с «.».

\makebibliography
% ИНСТРУКЦИЯ: список источников формируется автоматически из файла
% references.bib — в него нужно добавить все источники, на которые есть
% \cite в тексте (как минимум: книга по алгоритмам поиска и стандарт/
% документация выбранного языка). Источник без \cite в тексте в список
% не попадёт, и наоборот: \cite без записи в .bib даст ошибку.

\appendix
\chapter*{Приложение А}
\addcontentsline{toc}{chapter}{Приложение А}

% ИНСТРУКЦИЯ: здесь размещается полный исходный код программы. Каждый
% файл вставляется отдельной командой \lstinputlisting с параметрами:
%   language=  — язык (C++, Python, ...; для Swift может понадобиться
%                собственное описание языка или language={}),
%   caption=   — подпись, например {Заголовочный файл search.hpp},
%   label=     — метка, на которую ссылается раздел «Реализация
%                алгоритмов» (lst:search-hpp, lst:search-cpp,
%                lst:main-cpp, lst:test-cpp),
% и путём к файлу в фигурных скобках.
% Порядок листингов:
%   1) заголовочный файл с объявлениями функций поиска;
%   2) файл с реализацией функций поиска;
%   3) файл с главной функцией и меню;
%   4) файл с тестами.
% Проверьте, что все ссылки Листинг~\ref{...} из технологической части
% ведут на существующие метки (иначе в тексте будет «??»).

\end{document}